\chapter{Elliptische Kurven Kryptographie}
Nun wie wird aus abstrakter Mathematik ein kryptographisches Verfahren. In Kapitel 2 wurde schon Public-Key Kryptographie und das DL Problem betrachtet. Das selbe Prinzip soll nun genutzt werden, um ein kryptographisches Verfahren basierenden auf elliptischen Kurven zu implementieren. 
\label{sec:kryptographie}
\section{Diffie-Hellman}
\label{sec:diffie-hellman}
basiert auf \cite{washington08}, S.170-172
Wir betrachten eine möglichkeit einen geheimen Schlüssel zu etablieren nach Diffie und Hellman. Außerdem verwendet \cite{fuß25}, S. 195
\begin{enumerate}
    \item Alice und Bob einigen sich auf eine elliptische Kurve $E:y^2=x^3+Ax+B\mod p$ über einem endlichen Körper $\mathbb{F}_p$, sodass das diskrete Logarithmusproblem in $E(\mathbb{F}_p)$ schwer zu lösen ist.
    \item Alice wählt ein geheimes $a$, das ein Integer sein soll, berechne $P_a=aP$ und schickt $P_a$ zu Bob.
    \item BoB wählt einen geheimen Integer b, er berechnet $P_b=bP$ und schickt Alice $P_b$.
    \item Alice berechnet $aP_b=abP$
    \item Bob berechnet $bP_a=baP$
    \item Alice und Bob benutzen eine öffentliche Methode um einen Schlüssel von $abP$ zu extrahieren. Zum Beispiel könnten sie die letzten 256 bits der x-Koordinate von $abP$ als Schlüssel verwenden oder sie könnten eine Hash-Funktion auf der x-Koordinate definieren.
\end{enumerate}
Die einzige Information die ein Außenstehender sieht, wäre die Kurve $E$, den Körper $F_q$ und die Punkte $P, aP,bP$. Daher muss sie das \textbf{Diffie-Hellman Problem} lösen.

\begin{definition}[Diffie-Hellman Problem]
    Sei $P,aP,bP\in E(\mathbb{F}_q)$ gegeben berechne
    \[
        abP.
    \]
\end{definition}


\begin{beispiel}
    Wir betrachten die elliptische Kurve
    \[
        E\colon y^2 \equiv x^3+2x+2 \pmod{17}
    \]
    über dem Primkörper \(\mathbb F_{17}\). Durch explizite Bestimmung der Punkte erhält man
    \[
        \#E(\mathbb F_{17})=19.
    \]
    Da die Gruppenordnung \(19\) eine Primzahl ist, erzeugt jeder
    Punkt aus \(E(\mathbb F_{17})\setminus\{\mathcal O\}\) die gesamte
    Gruppe. Im Folgenden wird der Generator
    \[
        P=(0,6)
    \]
    gewählt. Die Berechnung der Ordnung, des Generators und alle weiteren Punktberechnungen wurde mit Hilfe des im Anhang~\ref{app:pythonprogramm} angegebene und selbstentickwelten Python-Programm durchgeführt.
    Alice wählt die geheime Zahl $a=2$ und berechnet den öffentlichen Schlüssel $A=aP=2P=(9,1)$, anschließend schickt Alice Bob $P_a$. Bob wählt $b=5$ und berechnet $B=bP=5P=(10,11)$, anschließend schickt Bob Alice $P_b$. Alice berechnet nun ihren gemeinsamen Schlüssel $S_A=aP_b=2\cdot(10,11)=(16,4)$, genauso wie Bob $S_B=bP_a=5\cdot(9,1)=(16,4)$.  Im Beispiel gilt konkret
    \[
        2B=2\cdot5P=10P=5\cdot2P=5A=(16,4).
    \]
    Öffentlich bekannt sind die Kurve \(E\), der Körper
    \(\mathbb F_{17}\), der Basispunkt \(P\) sowie die Punkte
    \(A=aP\) und \(B=bP\). Die Zahlen \(a\) und \(b\) sowie der
    gemeinsame Punkt \(S=abP\) werden dagegen nicht übertragen.

    Ein Angreifer müsste dann wie beschrieben, dass Diffie-Hellman-Problem lösen und \(abP\) finden. Eine Möglichkeit wäre \(a\) aus \(aP\) oder \(b\) aus \(bP\) mit dem diskreten Logarithmus Problem zu lösen. Anschließend könnte der Angreifer den gemeinsamen Punkt
    durch
    \[
        aB=a(bP)=abP
    \].Bei der hier betrachteten Gruppe mit nur \(19\) Elementen kann der
    diskrete Logarithmus durch einfaches Berechnen der Vielfachen
    \[
        P,2P,3P,\ldots
    \]
    gefunden werden. Bei kryptographisch verwendeten Kurven arbeitet man
    dagegen in einer Untergruppe von sehr großer Primzahlordnung. Eine
    vollständige Aufzählung der Vielfachen von \(P\) ist dort praktisch
    nicht durchführbar.
\end{beispiel}

Um eine sichere Konfiguration für Diffie-Hellman zu erhalten muss $p$ relativ groß ($p>2^{160})$ und daher muss man eine große Anzahl von Vielfachen von einem Punkt $P$ berechnen
\[
    mP = P +P +P+\dots P\in E(\mathbb{F}_p).
\]
Wenn man
\cite{silverman06}, S. 30


Diffie-Hellman in dieser Form führt nur dazu, dass ein gemeinsamer Schlüssel ausgetauscht wird, es gibt noch keine Beschreibung wie genau die Verschlüsslung und Entschlüsslung funktioniert. Dazu kann man zum Beispiel das ElGamal Public Key Verfahren verwenden. 
Alice möchte Bob eine Nachricht senden. Dafür muss Bob aber seinen öffentlichen Schlüssel etablieren. Dafür wählt er eine elliptische Kurve $E$ über einem endlichen Körper $\mathbb{F}_q$, so dass das DL Problem schwer zu lösen ist auf $E\mathbb{F}_q$ (wie genau diese Bedingungen sein muss wird in Abschnitt \ref{sicherheitsbetrachtung} betrachtet). Außerdem wählt er auch einen Punkt $P$ auf $E$. Er wählt wieder einen geheimes $b\in\mathbb{Z}$ und berechnet $B=bP$. Die elliptische Kurve $E$, der endliche Körper $\mathbb{F}_q$ und die Punkte $P$ und $B$ macht er öffentlicht und dienen ihm als öffentliche Schlüssel. Wenn Alice Bob eine Nachricht schicken möchte geht sie wie folgt vor
\begin{enumerate}
    \item Alice lädt Bobs öffentlichen Schlüssel runter
    \item Sie schreibt ihre Nachricht als ein Punkt $M\in E(\mathbb{F}_q)$
    \item Nun wählt Alice einen geheimen Schlüssel $k$ und berechnet $M_1=kP$
    \item Anschließend berechnet Alice $M_2=M+kB$
    \item Nun sendet Alice $M_1,M_2$ an Bob
\end{enumerate}

Bob entschlüsselt die Nachricht indem er \(M=M_2-bM_1\) berechnet. Das funktioniert da Bob $b$ kennt und daher 
\[
    M_2-bM_1=(M+kB)-s(kP)= M+k(bP)-bkP=M
\]
berechnen kann. Eine dritte Person Eve die die Nachrichten $M_1$ und $M_2$ abfängt, kann nur durch Lösen des DL-Problems $b$ aus den öffentlichen Informationen $P$ und $B$ berechnen. Wenn sie $b$ nicht errechnen, kann exestiert kein Weg um die Nachricht $M$ zu entschlüsseln.
\begin{bemerkung}
    In der Regel wird mit asymmetrischen Verfahren nur ein symmetrischer Schlüssel ausgetauscht und daher handelt es sich meist bei der Nachricht um einen Punkt auf der Kurve der als Geheimnis dient zum Beispiel wird die x-Koordinate des gemeinsamen Schlüssels verwendet und mit einer Hashfunktion verschlüsselt und dann als symmetrischer Schlüssel verwendet (vgl. \cite{guneysu24}, S.291). Wenn tatsächlich eine richtige Nachricht wie zum Beispiel "Hallo" geschickt werden soll, dann muss diese zuerst in ihre Bitdarstellung überführt werden, zum Beispiel mit Hilfe von ASCII und anschließend dann mit einer Abbildung auf einen Punkt auf der Kurve abgebildet werden.
\end{bemerkung}

\begin{bemerkung}
    Wenn Alice das selbe $k$ zweimal verwendet, ist $M_1'=kP=M_1$. Eve kann das Wissen nutzen, um $M$ zu berechnen indem sie $M_2'-M_2= M'-bM_1'-(M-bM_1)=M'-M$ berechnet. Wenn Eve nun Wissen über $M$ erhält, kann sie dadurch auch automatisch $M'$ berechnen und alle weiteren Nachrichten die anschließend mit $k$ verschlüsselt werden.
\end{bemerkung}


\cite{washington08}, S. 174-176
\begin{definition}[ElGamal Verschlüsslung]

    Es exestiert nun zum erzeugten Schlüssel $(ab)\cdot G$ eine Funktion
    \[
    F: E \to E, P\mapsto P+(ab)\cdot G
    \]
    Es gibt dann eine Verschlüsslungsfunktion
    \[
    P\mapsto (B,Q)= (b\cdot G,P+b\cdot A)
    \]
    und eine Entschlüsslungsfunktion
    \[
    (B,Q)\mapsto Q-a\cdot B=Q-(ab)\cdot G=Q-b\cdot A=P
    \]
\end{definition}
Exklusives Geheimnis von Alice ist einzig die Zahl a. Bobs Geheimnis ist nur die Zahl b. Beide Zahlen können aus den Punkten $A=a\cdot G$ und $B=b\cdot G$ nicht zurückgerechnet werden.
\begin{bemerkung}
    Wie beim klassischen Elgamal-Verfahren muss Bob jedes mal eine neue Zahl $b$ auswählen. Denn ansonsten kann aus dem Geheimtext \(b\cdot G, P+b\cdot A\) mit Klartext $P$, der Kurvenpunkt $b\cdot A$ bestimmt werden. Dadurch können alle andere Punkte die mit $b$ verschlüsselt wurden entschlüsselt werden.
\end{bemerkung}
\section{ECDSA}

\section{Eventuell noch ECIES \cite{washington08},S.180ff.}
\section{Warum eliptische Kurven verwenden?}
Der beste Algorithm um ECDLP zu lösen ist exponential, darum werden Elliptische Kurven Gruppen verwendet. (An Introduction to the Theory of Elliptic Curves, S.6)
Genauer der beste Algorithmus um ECDLP zu lösen über $\mathbb{F}_p$ braucht $O(\sqrt p)$ Zeit.


Ephemeral-Diffie-Hellman-Key-Agreement mit elliptischen
Kurven

ECDSA


\section{Sicherheitbetrachtung und Wahl von sinnvollen Bedingugen}
\label{sicherheitsbetrachtung}
\textbf{Warum elliptische Kurven?}
Praktisch stellt sich das RSA-Verfahren erst ab einer Schlüsselänge von 3000 Bit als sicher heraus. Entsprechend aufwendig gestaltet sich dort das Ver- und
Entschlüsseln. Um das Verschlüsseln effizient zu gestalten, kann ein kleiner Exponent e gewählt werden. Elliptische Kurven haben sich in der modernen Kryptographie etabliert, weil sie bei vergleichbarem Sicherheitsniveau deutlich kürzere Schlüssellängen erfordern als klassische Verfahren wie RSA.     So bietet ein RSA-Schlüssel mit 4096 Bit ein Sicherheitsniveau, das mit Verfahren auf Basis elliptischer Kurven bereits mit einer Schlüssellänge von etwa 313 Bit erreicht werden kann (vgl. \cite{washington08}, S. 169). Dies führt zu geringerem Speicherbedarf, ermöglicht kompaktere Hardware und reduziert den Energieverbrauch.
Elliptische Kurven bieten nun eine Möglichkeit um, wie wir in den folgenden Kapiteln sehen werden, ein schwer zu lösendes diskretes Logarithmus Problem zu stellen.
Faktorisierung großer Zahlen (Computeralgebra Kapitel 5, ab S.167)
\textbf{Weil-Pairing?}


\begin{itemize}
    \item Punktzählung
    \item Baby Step Giant Step
    \item Pohlig-Hellman zeigt, dass man Untergruppen großer Primzahlordnung braucht
    \item MOV 
\end{itemize}
\section{Performancesteigerung und Vergleich von ECC zu anderen Verfahren}

Performance wäre zu finden bei \cite{fuß25}, S.201

\cite{fuß25}, S.199-200

\begin{figure}[htbp]
    \centering
    \includegraphics[width=\linewidth]{empfehlung_der_nist.png}
    \caption{Empfehlungen des NIST, entnommen aus \cite[S.~219]{buell24}.}
    \label{fig:nist-empfehlungen}
\end{figure}

\section{Angriffe auf ECC}
In Kapitel 2 \ref{baby-step-giant-step} wurde der Baby Step Giant Step Algorithmus verwendet nun soll dieser nochmal an einem Beispiel für elliptische Kurven aufgegriffen werden.
\begin{beispiel}
    Es wird erneut das Beispiel aus Kapitel 4 aufgegriffen, aber es wird angenommen, dass die Gruppenordnung unbekannt sei. Sei $G=E(\mathbb{F}_{17})$, wobei $E$ gegeben ist durch $y^2=x^3+2x+2$. Sei $P=(0,6)$ und $Q=(3, 16)=kP$. nach dem Satz von Hasse \ref{def:hasse-schranke} kann eine obere Schranke ausgerechnet werden
    \[ 
    |\#E(\mathbb{F}_{17})-17-1 |\leq 2\sqrt 17
    \],
    also die Ordnung maximal 26. Daher sei $m=\lceil \sqrt{26}\rceil=6$. Die Punkte $iP$ für $1\leq i< 6$ sind
    \[
    (0, 6), (9, 1), (6, 3), (7, 6), (10, 11)
    \]
    Für die linke Seite wird dann $Q-jmP$ berechnet für $j=0,1,2,3,4$
    \[
    (3,16),(13, 10),(0, 6),(10, 6),(5,16),(9,1).
    \]
    Es wird bei $j=2$ beendet, da der Punkt mit $1P$ übereinstimmt. Nun kann k einfach nachgerechnet werden, 
    \[
    (30,40)=kP=(1+2\cdot6)P= 13P
    \]
    und damit erhalten wir $k=13$ als das gesuchte Ergebnis.
    Alles wurde mit Hilfe von \ref{app:pythonprogramm}, die Konsolenausgabe ist in \ref{fig:python-2}.
\end{beispiel}
Nun wird das Pohlig Hellmann Verfahren aus Kapitel 2 \ref{pohlig-hellman} ebenefalls für elliptische Kurven aufgegriffen.
\begin{beispiel}
    Sei $G=E(\mathbb{F}_{599})$ für die elliptische Kurve
    $E\colon y^2=x^3+1$. Gegeben seien die Punkte
    $P=(60,19)$ und $Q=(277,239)$. Das Python-Programm
    liefert $|G|=\operatorname{ord}(P)=600$. Gesucht ist
    ein $k\in\{0,\dots,599\}$ mit $Q=kP$.
    Wegen $600=2^3\cdot3\cdot5^2$ bestimmen wir wie in
    Kapitel~2 zunächst $k$ modulo $8$, $3$ und $25$.
    Alle Koordinatenrechnungen erfolgen modulo $599$.

    \smallskip
    \noindent\textbf{Bestimmung von $k\bmod 8$.}
    Wir schreiben $k\equiv z_0+2z_1+4z_2\pmod 8$
    mit $z_i\in\{0,1\}$. Der Punkt
    $R=\frac{600}{2}P=300P=(598,0)$ hat die Ordnung $2$.
    Die Berechnung ergibt $300Q=e_G=0R$, also $z_0=0$.
    Die weiteren Ziffern ergeben sich aus
    \[
    \begin{aligned}
        Q_1&=\frac{600}{2^2}(Q-z_0P)
             =150Q=(598,0)=R,\\
        Q_2&=\frac{600}{2^3}(Q-(z_0+2z_1)P)
             =75(Q-2P)=75(35,243)=e_G.
    \end{aligned}
    \]
    Somit sind $z_1=1$ und $z_2=0$, weshalb
    $k\equiv 0+2\cdot1+4\cdot0=2\pmod 8$ gilt.

    \smallskip
    \noindent\textbf{Bestimmung von $k\bmod 3$.}
    Hier ist $R=\frac{600}{3}P=200P=(0,1)$.
    Aus $200Q=(0,598)=2R$ folgt unmittelbar
    $k\equiv2\pmod 3$.

    \smallskip
    \noindent\textbf{Bestimmung von $k\bmod 25$.}
    Wir schreiben $k\equiv z_0+5z_1\pmod{25}$
    mit $z_0,z_1\in\{0,\dots,4\}$ und setzen
    $R=\frac{600}{5}P=120P=(84,179)$.
    Aus $120Q=(84,179)=R$ folgt $z_0=1$.
    Anschließend berechnen wir
    \[
        Q_1=\frac{600}{5^2}(Q-P)
            =24(130,129)=(491,465)=3R.
    \]
    Daher ist $z_1=3$ und damit
    $k\equiv1+5\cdot3=16\pmod{25}$.

    \smallskip
    \noindent\textbf{Zusammensetzen der Ergebnisse.}
    Aus $k\equiv2\pmod 8$ und $k\equiv2\pmod 3$
    folgt $k\equiv2\pmod{24}$.
    Mit $k=2+24t$ ergibt die verbleibende Kongruenz
    \[
        2+24t\equiv16\pmod{25}
        \quad\Longleftrightarrow\quad
        t\equiv11\pmod{25}.
    \]
    Somit gilt $k\equiv266\pmod{600}$.
    Der gesuchte diskrete Logarithmus ist also $k=266$.
\end{beispiel}

\subsection{MOV-Algorithmus}
\cite{washington08} S. 156--157
\begin{lemma}
    Sei $E$ eine elliptische Kurve über dem Körper $\mathbb{F}_q$ mit
    \[
        \operatorname{Spur}(\phi_q)
        = q+1-\#E(\mathbb{F}_q)=0.
    \]
    Sei $n\in\mathbb{Z}_{>0}$. Existiert ein Punkt
    $P\in E(\mathbb{F}_q)$ der Ordnung $n$, so gilt
    \[
        E[n]\subseteq E(\mathbb{F}_{q^2}).
    \]
\end{lemma}

Dadurch kann das diskrete Logarithmusproblem auf einer supersingulären elliptischen Kurve mit $\operatorname{Spur}(\phi_q)=0$ auf ein diskretes Logarithmusproblem in der multiplikativen Gruppe $\mathbb{F}_{q^2}^{\times}$ reduziert werden. Da hierfür lediglich eine
Körpererweiterung vom Grad $2$ benötigt wird, können effizientere Verfahren zur Berechnung diskreter Logarithmen in endlichen Körpern
angewendet werden.

Auch für supersinguläre Kurven mit $\operatorname{Spur}(\phi_q)\neq 0$ ist eine solche Reduktion möglich. Dabei kann ein rweiterungsgrad $m\in\{3,4,6\}$ erforderlich sein. Diese Erweiterungsgrade sind weiterhin klein und ermöglichen daher eine Beschleunigung der Berechnung des diskreten Logarithmus.